\subsection{SA19 H1--H2: mathematical identification and the chronology gate}

\paragraph{H1: component-group duality.}
The boxed statement on 43:117 identifies the component groups of dual
abelian varieties as dual with values in $\QQ/\ZZ$. This matches
SGA~7~I, Expos\'e~IX, (1.2.1) and Conjecture~1.3, printed p.~323
(PDF p.~328 in the opened original scan): the pairing is constructed,
while its perfectness is stated separately as a conjecture
\cite{SGA7I}. The original page was visually checked.
For a complete discrete valuation field with perfect residue field of
characteristic $p>0$, perfectness is now Suzuki's Theorem~C,
p.~4 of the opened author version; see its setup on p.~2 and
Theorems~A--B on pp.~3--4 \cite{SuzukiPairing2020}.
An algebraically closed residue field satisfies this hypothesis.
The derived duality and descent in those theorems give a substantive
modern counterpart to the programme on 43:116--120. This is a
known theorem, not a new proof here. The residue-field hypothesis
cannot simply be dropped: Bosch--Lorenzini's Lemma~6.1 and
Example~6.2, pp.~393--395, give Jacobians over imperfect residue
fields for which the pairing has nonzero radical
\cite{BoschLorenzini2002}; these examples are outside the manuscript's
algebraically closed setting.

\paragraph{H2: three-variable Ext and biextensions.}
The boxed three-variable Ext on 43:16 is labelled the Lang pairing.
On 43:21 the tensor product of two resolutions, the first cohomology
of a Hom complex, and the equivariant cup product give additional
evidence for a biextension interpretation. Some decorations and
inputs remain uncertain in the transcription; three-variable
notation alone does not specify an exact category or a construction.
The precise published comparison is
\[
 \operatorname{Biext}^{1}(P,Q;G)
   \simeq \operatorname{Ext}^{1}(P\mathbin{\otimes}^{\mathbf L}Q,G),
\]
SGA~7~I, Expos\'e~VII, Corollary~3.6.5, printed p.~205
(PDF p.~210), following the categorical construction in 3.6.4,
p.~204 \cite{SGA7I}. The original formula was visually checked;
the corollary states degree one, not an unrestricted all-degree
comparison.

Mumford's \emph{Bi-extensions of formal groups}, \S\S2--3,
original pp.~310--314, gives the torsor definition with two compatible
extension laws and then the formulation for group functors
\cite{Mumford1969}. The official TIFR retypeset volume retains
original page numbers in its margins; its running article pages
321--337 differ from the original 307--322. Expos\'e~VII's introduction,
printed p.~133, explicitly attributes the introduction of
biextensions in the formal-group setting to Mumford. Any historical
priority claim must address that contemporaneous attribution.

\paragraph{The dating inference is invalid; priority remains Undetermined.}
The dates cited in SA19 H1--H2 do not date the relevant scientific
leaves. On 43:144, June~1962 dates Shatz's cited Harvard thesis.
On 43:87, May~1964 dates Lang's cited Bourbaki talk; that bibliography
also cites N\'eron's 1964 paper. These are bibliographic dates,
which do not give an upper bound for the dates of other leaves.
On 43:118, September~1963 occurs in a cancelled administrative list
mentioning James and Bott. Adjacency to 43:117 does not establish
when the latter's scientific content was written. The supplied
transcriptions of 43:16, 43:21 and 43:117 carry no explicit dates,
and the inspected cover sheets provide none.

The official Montpellier catalogue gives cote~43 the tentative span
\emph{[à partir de 1961-vers 1968]}, marking the notes themselves
undated. Footnote~37 records a verso with Bourbaki exposé~338
(1967--1968); the catalogue also warns that versos are often
unrelated to the mathematical notes \cite{MontpellierCatalogue}.
This does not support a dossier-wide upper bound of 1964 or 1965.
Nor does that late verso date the particular scientific leaves.
The catalogue dates cote~44 separately to 1964; this cannot be
transferred to cote~43.

Thus the mathematical identifications are supported, but the offered
1962--64 chronology does not prove pre-SGA-7 priority. This refutes
the proposed inference, not the possibility of early composition.
The exact remaining gate is independently secure dating or provenance
for the particular scientific leaves, together with an argument
connecting that evidence to those texts. A dated quotation or
independently dated matching draft could close the gate; a repeated
dossier-wide date range cannot.
